\begin{table}[ht]
\caption{\textbf{TheoremQA Theorems example.} The query uses the Pigeonhole Principle, which is the gold document from ProofWiki. The negative document is retrieved by BM25, which is not relevant to the question}
\centering
% \resizebox{\textwidth}{!}{
\begin{tabular}{p{13cm}}
\toprule
% \hspace{5.5cm} \sustainable \\
% \midrule
\textit{\textbf{Query}} \\
\midrule
Mary is planning to bake exactly 10 cookies, and each cookie may be one of three different shapes -- triangle, circle, and square. Mary wants the cookie shapes to be a diverse as possible. What is the smallest possible count for the most common shape across the ten cookies?\\
\midrule
\textit{\textbf{Chain-of-thought reasoning to find documents}} \\
\midrule
This problem is similar to the formulation of Pigeonhole Principle. We may look into the principle and find hints about how this problem may be solve using that principle.\\
\midrule
\textit{\textbf{Example positive document}} \\
\midrule
Pigeonhole Principle Tags: Pigeonhole Principle, Named Theorems, Combinatorics  
\\begin\{theorem\} Let $S$ be a finite set whose cardinality is $n$. Let $S_1, S_2, \ldots, S_k$ be a partition of $S$ into $k$ subsets. Then: :at least one subset $S_i$ of $S$ contains at least $\ceiling{\dfrac n k}$ elements where $\ceiling{\, \cdot \,}$ denotes the ceiling function. \\end\{theorem\}  
\\begin\{proof\} {{AimForCont}} no subset $S_i$ of $S$ has as many as $\ceiling {\dfrac n k}$ elements. 
Then the maximum number of elements of any $S_i$ would be $\ceiling {\dfrac n k} - 1$. 
So the total number of elements of $S$ would be no more than $k \left(\ceiling {\dfrac n k} - 1\right) = k \ceiling {\dfrac n k} - k$. There are two cases: :$n$ is divisible by $k$ :$n$ is not divisible by $k$. 
Suppose $k \backslash n$. Then $\ceiling {\dfrac n k} = \dfrac n k$ is an integer and: :$k \ceiling {\dfrac n k} - k = n - k$ Thus: :$\card S = \sum_{i \mathop = 1}^k \card {S_i} \le n - k < n$ This contradicts the fact that $\card S = n$. 
Hence our assumption that no subset $S_i$ of $S$ has as many as $\ceiling {\dfrac n k}$ elements was false. 
Next, suppose that $k mid n$. Then: :$\card S = k \ceiling{\dfrac n k} - k < \dfrac {k (n + k) } k - k = n$ and again this contradicts the fact that $\card S = n$. 
In the same way, our assumption that no subset $S_i$ of $S$ has as many as $\ceiling {\dfrac n k}$ elements was false. Hence, by Proof by Contradiction, there has to be at least $\ceiling {\dfrac n k}$ elements in at least one $S_i \subseteq S$. {{qed}} \\end\{proof\}  \\
\midrule
\textit{\textbf{Example negative document}} \\
\midrule
\\begin\{definition\}[Definition:Conditional/Semantics of Conditional]Let $p \implies q$ where $\implies$ denotes the conditional operator.$p \implies q$ can be stated thus:* '''''If'' $p$ is true ''then'' $q$ is true.'''* '''$q$ is true ''if'' $p$ is true.'''* '''(The truth of) $p$ ''implies'' (the truth of) $q$.'''* '''(The truth of) $q$ ''is implied by'' (the truth of) $p$.'''* '''$q$ ''follows from'' $p$.'''* '''$p$ is true ''only if'' $q$ is true.'''The latter one may need some explanation. $p$ can be either true or false, as can $q$. But if $q$ is false, and $p \implies q$, then $p$ can not be true. Therefore, $p$ can be true ''only if'' $q$ is also true, which leads us to our assertion.* '''$p$ is true ''therefore'' $q$ is true.'''* '''$p$ is true ''entails'' that $q$ is true.'''* '''$q$ is true ''because'' $p$ is true.'''* '''$p$ ''may'' be true ''unless'' $q$ is false.'''* '''''Given that'' $p$ is true, $q$ is true.'''* '''$q$ is true ''whenever'' $p$ is true.'''* '''$q$ is true ''provided that'' $p$ is true.'''* '''$q$ is true ''in case'' $p$ is true.'''* '''$q$ is true ''assuming that'' $p$ is true.'''* '''$q$ is true ''on the condition that'' $p$ is true.'''Further colloquial interpretations can often be found in natural language whose meaning can be reduced down '''$p$ only if $q$''', for example:* '''$p$ is true ''as long as'' $q$ is true'''::Example::::''"Mummy, can I go to the pictures?"'':::''"'''As long as''' you've done your homework. Have you done your homework? No? Then you cannot go to the pictures."''::In other words::::''"You can go to the pictures '''only if''' you have done your homework."''::Using the full language of logic::::''"If it is true that you are going to the pictures, it is true that you must have done your homework."''* '''$p$ is true ''as soon as'' $q$ is true''':::''"Are we going to this party, then?"'':::''"'''As soon as''' I've finished putting on my makeup.''"::The analysis is the same as for the above example of '''as long as'''.\\end\{definition\}
\\
\bottomrule
\end{tabular}
% }
\label{tab:example_theoremqa_theorems}
\end{table}